\section{Conclusion}
\label{sec:conclusion}

Big Data platforms remain expensive to provision and difficult to reproduce,
which narrows access for academic units and smaller organisations. The wider
literature increasingly emphasises Infrastructure-as-Code and scripted
deployment as safeguards against configuration drift; situating a full Hadoop
Data Lake on Oracle Cloud Infrastructure Free Tier ARM instances extends that
line of work to a deliberately minimal budget and hardware envelope where few
published end-to-end references exist.

This paper presented an open-source pipeline that provisions a four-node ODP
cluster through Terraform, Ansible, and the Ambari API, with automated recovery
on failed installs. Installation profiles let the same codebase target
different workloads; the Star Schema Benchmark was executed on the
data-science profile to stress Spark and Hive under tight memory and
CPU ceilings.

Empirically, Spark reduced mean query time relative to Hive on every SSB query
by \PaperHSRedPctMin\% to \PaperHSRedPctMax\% per query. In-memory caching did
not deliver a uniform gain: lighter queries slowed by roughly
\PaperCacheOverheadNarrativeMinPct\% to \PaperCacheOverheadNarrativeMaxPct\%
against uncached Spark because the 9.3~GB lineorder table could not
fit entirely in RAM, whereas several heavier multi-join queries benefited from
partial dimension caching (up to about \PaperQThreeTwoCacheImprovementPct\%
improvement on Q3.2). Averaged over all thirteen queries, the cached
configuration remained \PaperCacheVsSparkMeanPenaltyPct\% slower than uncached
Spark. All benchmark runs completed without unrecoverable failures,
so the stack proved usable for sustained testing despite the constraints.

The study has clear limits. It used a fixed four-node cluster with no x86
baseline, so architectural effects cannot be separated from the ARM
free-tier envelope. Four OCPUs and 24~GB of RAM cap parallelism and cacheable volume,
and differences in JVM behaviour or packaging between ARM and x86 should be
weighed before any production adoption.

Future work could exercise streaming-oriented paths with Spark Structured Streaming, adopt
columnar formats such as ORC or Parquet to cut disk I/O, run a controlled ARM
versus x86 comparison, and explore horizontal scale through additional
tenancies or paid instances with more executors to relax the bottlenecks seen
here.
